%! Setting Up the SVD — Unit 7, Lesson 7.0 — Warm-Up (Answer Key)
\documentclass[10pt]{article}
\usepackage{linalg-article}
\usepackage{linalg-key}   % pulls in -boxes; do NOT also load -boxes
\newcommand{\TallMath}[1]{$\displaystyle #1\rule[-1.4em]{0pt}{3.2em}$}

\begin{document}

\pageheader{Unit 7, Lesson 7.0}{Warm-Up --- Answer Key}
\namedateperiod

\vspace{0.12in}

\begin{notesbox}{Getting Ready: Ideas We'll Use Today}
\small These three ideas from Units~1 and~4 set up today's picture of the SVD. Answer quickly.

\vspace{4pt}
\textbf{1. Matrix times a vector (Unit 1).} Compute $A\mathbf{v}$ for
$A=\begin{bmatrix} 2 & 1 \\ 1 & 2 \end{bmatrix}$ and $\mathbf{v}=\begin{bmatrix} 1 \\ 1 \end{bmatrix}$:
\[
  A\mathbf{v}=\begin{bmatrix} 2 & 1 \\ 1 & 2 \end{bmatrix}\begin{bmatrix} 1 \\ 1 \end{bmatrix}
  = \begin{bmatrix}{\color{keyred}\mathbf{3}}\\[2pt]{\color{keyred}\mathbf{3}}\end{bmatrix}.
\]

\vspace{4pt}
\textbf{2. Perpendicular vectors (Unit 4).} Two vectors are \emph{perpendicular} when their dot
product is $0$. Are $\begin{bmatrix} 1 \\ 1 \end{bmatrix}$ and $\begin{bmatrix} 1 \\ -1 \end{bmatrix}$
perpendicular?
\[
  \begin{bmatrix} 1 \\ 1 \end{bmatrix}\cdot\begin{bmatrix} 1 \\ -1 \end{bmatrix}
  = (1)(1)+(1)(-1) = {\color{keyred}\mathbf{0}}\quad\Rightarrow\quad \text{perpendicular? }\ans{yes}
\]

\vspace{4pt}
\textbf{3. Length and unit vectors (Unit 4).} The length of $\begin{bmatrix} 1 \\ 1 \end{bmatrix}$ is
$\sqrt{1^2+1^2}={\color{keyred}\mathbf{\sqrt{2}}}$. To make it a \emph{unit} vector (length $1$), divide by its length:
\[
  \frac{1}{{\color{keyred}\mathbf{\sqrt{2}}}}\begin{bmatrix} 1 \\ 1 \end{bmatrix}.
\]
\end{notesbox}

\begin{teachernote}
All three skills reappear today. Item~1 is a matrix stretching a vector (the $\begin{bsmallmatrix}1\\1\end{bsmallmatrix}$
direction goes to $\begin{bsmallmatrix}3\\3\end{bsmallmatrix}$ --- three times as long). Items~2--3 are the
Unit~4 tools we need to describe the SVD's \emph{perpendicular} input and output directions and to
report them as \emph{unit} vectors. Debrief item~2 aloud: a dot product of $0$ means a right angle.
\end{teachernote}

\end{document}
